O trabalho partiu de 1.159.764 faixas e descreveu as diferenças de distribuição por gênero, período e atributos acústicos. A estratificação por gênero e as metas de precisão determinaram um sorteio de 9.526 faixas. Os 3.869 vínculos aceitos desse sorteio demonstraram cobertura seletiva, especialmente por gênero e popularidade. A avaliação final usou outra seleção, de 9.526 faixas com letra, feita em um quadro consolidado de 80.943 vínculos. Essa distinção permite usar o primeiro sorteio para caracterizar a recuperação e a coorte vinculada para comparar os modelos.

A regressão ordinal com descritores tradicionais teve acurácia balanceada de 56,24\%, próxima aos 56,14\% do Baseline. Com embeddings, houve maior sensibilidade da classe alta ao custo de menor precisão nessa classe e menor sensibilidade nas demais. Os intervalos condicionados às previsões incluíram zero em todos os contrastes da métrica principal. O conjunto dos resultados indica que a letra alterou a distribuição dos erros, com ganho global limitado nas condições avaliadas.

A comparação sob o mesmo desenho de validação quantifica o desempenho incremental de cada conjunto de preditores. O alcance continua delimitado pela cobertura textual, pelo escore disponível e pelas faixas avaliadas. Uma replicação com faixas de outro período e cobertura caracterizada serviria para verificar a estabilidade dos resultados fora desse domínio.

